\documentclass[../../../include/open-logic-section]{subfiles}

\begin{document}

\olfileid[fr]{sth}{card-arithmetic}{ch}

\olsection{L'hypothèse du continu}

Le résultat précédent suggère, à juste titre,
que l'exponentiation cardinale serait assez
\emph{simple} si les cardinaux infinis se comportaient
de façon «~régulière~» vis-à-vis des puissances de $2$,
c'est-à-dire, d'après le
\olref[opps]{lem:SizePowerset2Exp}, vis-à-vis
de l'opération qui forme l'ensemble des parties.
Nous ne pouvons toutefois pas supposer qu'ils
se comportent \emph{effectivement} ainsi.

Pour commencer à éclaircir la question,
introduisons une notation commode.

\begin{defn}
Si $\cardsucc{\cardfont{a}}$ désigne le plus
petit cardinal strictement supérieur à
$\cardfont{a}$, nous définissons deux suites
transfinies :
\begin{align*}
	\aleph_{0} &\defis \omega &
	\beth_{0} &\defis \omega\\
	\aleph_{\alpha \ordplus 1}
	&\defis \cardsucc{(\aleph_{\alpha})} &
	\beth_{\alpha+1}
	&\defis \cardexpo{2}{\beth_{\alpha}}\\
	\aleph_{\alpha}
	&\defis \bigcup_{\beta< \alpha} \aleph_{\beta} &
	\beth_{\alpha}
	&\defis \bigcup_{\beta < \alpha}\beth_{\beta}
	& \text{si $\alpha$ est un ordinal limite}.
\end{align*}
\end{defn}

La définition de $\cardsucc{\cardfont{a}}$
est légitime : d'après le
\olref[cardinals][classing]{lem:NoLargestCardinal},
pour chaque cardinal $\cardfont{a}$,
il existe un cardinal qui lui est supérieur ;
la récurrence transfinie garantit qu'il en
existe un \emph{plus petit}.




Le reste de la définition de ces deux suites
est donné par récursion transfinie.

Cantor a introduit la notation «~$\aleph$~».
Il s'agit de l'\emph{aleph}, première lettre
de l'alphabet hébreu et du mot hébreu signifiant
«~infini~». Peirce a introduit la notation
«~$\beth$~» : le \emph{beth} est la deuxième
lettre de l'alphabet hébreu.\footnote{
Peirce a employé cette notation dans une lettre
à Cantor datée de décembre 1900. Malheureusement,
il y donnait aussi un mauvais argument selon
lequel $\beth_\alpha$ n'existerait pas
pour $\alpha \geq \omega$.}
Ces notations nous fournissent des cardinaux
infinis.

\begin{prop}
$\aleph_\alpha$ et $\beth_\alpha$ sont
des cardinaux pour tout ordinal $\alpha$.
\end{prop}

\begin{proof}
Les deux assertions se démontrent par une
récurrence transfinie simple.
$\aleph_0 = \beth_0 = \omega$ est un cardinal
par le \olref[cardinals][classing]{omegaisacardinal}.
Si $\aleph_\alpha$ et $\beth_\alpha$ sont des
cardinaux, alors $\aleph_{\alpha+1}$ et
$\beth_{\alpha+1}$ sont explicitement définis
comme des cardinaux. Enfin, l'union d'un ensemble
de cardinaux est un cardinal, d'après la
\olref[cardinals][classing]{unioncardinalscardinal}.
\end{proof}
\noindent
De plus, tout cardinal infini est un aleph :

\begin{prop}
Si $\cardfont{a}$ est un cardinal infini,
alors $\cardfont{a} = \aleph_\gamma$
pour un unique $\gamma$.
\end{prop}

\begin{proof}
Raisonnons par récurrence transfinie sur les
cardinaux infinis. Le cas initial est
$\omega=\aleph_0$. Pour l'étape de récurrence,
supposons que, pour tout cardinal infini
$\cardfont{b} < \cardfont{a}$, on ait
$\cardfont{b} = \aleph_{\gamma_\cardfont{b}}$.
Si $\cardfont{a} =
\cardsucc{\cardfont{b}}$ pour un tel
$\cardfont{b}$, alors
$\cardfont{a} =
\cardsucc{(\aleph_{\gamma_\cardfont{b}})}
= \aleph_{\gamma_\cardfont{b}+1}$.
Si $\cardfont{a}$ n'est le successeur d'aucun
cardinal, alors, puisque les cardinaux sont des
ordinaux,
$\cardfont{a}
= \bigcup_{\omega \leq \cardfont{b} < \cardfont{a}}
\cardfont{b}
= \bigcup_{\omega \leq \cardfont{b} < \cardfont{a}}
{\aleph_{\gamma_\cardfont{b}}}$ ;
donc $\cardfont{a} = \aleph_\gamma$, où
$\gamma =
\bigcup_{\omega \leq \cardfont{b} < \cardfont{a}}
\gamma_\cardfont{b}$.
%
%
\footnote{Note de l'édition française : l'original
omet le cas initial $\omega$, puis emploie
$\gamma_\cardfont{b}$ pour des cardinaux
éventuellement finis, auxquels l'hypothèse
de récurrence ne s'applique pas. La récurrence
est ici restreinte aux cardinaux infinis,
avec son cas initial.}
\end{proof}

Puisque tout cardinal infini est un aleph,
une question se pose : tout cardinal infini
est-il aussi un beth ? Si \emph{tel était}
le cas, les cardinaux infinis se comporteraient
effectivement de façon «~régulière~»
vis-à-vis des ensembles de parties.
Nous aurions alors :

\begin{defish}
\emph{Hypothèse généralisée du continu} (GCH).
$\aleph_\alpha = \beth_\alpha$ pour tout $\alpha$.
\end{defish}

De plus, si la GCH était vraie, nous pourrions
simplifier considérablement l'exponentiation
cardinale. En particulier, pour un cardinal
infini $\cardfont{a}$ et $0<\cardfont{b}<\cardfont{a}$,
nous pourrions montrer que
$\cardfont{a}
\leq \cardexpo{\cardfont{a}}{\cardfont{b}}
\leq \cardsucc{\cardfont{a}}$.\footnote{
Note de l'édition française : l'original énonce
cette borne sans préciser que l'exposant doit
être non nul ; pour $\cardfont{b}=0$,
$\cardexpo{\cardfont{a}}{0}=1$,
et la borne inférieure échoue. Le cardinal
$\cardfont{a}$ est ici infini.}
Nous pourrions ensuite donner des conditions
précises permettant de savoir laquelle des
deux possibilités se réalise, à savoir
$\cardfont{a}
= \cardexpo{\cardfont{a}}{\cardfont{b}}$
ou
$\cardexpo{\cardfont{a}}{\cardfont{b}}
= \cardsucc{\cardfont{a}}$.\footnote{
La condition dépend de la \emph{cofinalité}.}

Mais la GCH est une \emph{hypothèse}, et non
un \emph{théorème}. En fait, \citet{Godel1938}
a démontré que si $\ZFC$ est cohérente, alors
$\ZFC + \text{GCH}$ l'est aussi. On a découvert
par la suite que l'on pouvait tout autant ajouter
$\lnot$GCH à $\ZFC$. Considérons en effet le
plus simple cas \emph{non trivial} de la GCH :

\begin{defish}
\emph{Hypothèse du continu} (CH).
$\aleph_1 = \beth_1$.
\end{defish}

\citet{Cohen1963} a démontré que, si $\ZFC$
est cohérente, alors
$\ZFC + \lnot\text{CH}$ l'est aussi.
L'hypothèse du continu est donc indépendante
de $\ZFC$.

Elle tire son nom du fait que «~le continu~»
est un autre nom de la droite réelle, $\Real$.
Le \olref[opps]{continuumis2aleph0} nous dit que
$\card{\Real} = \beth_1$. L'hypothèse du continu
affirme donc qu'il n'existe aucun cardinal
strictement compris entre celui des nombres
naturels, $\aleph_0 = \beth_0$, et celui
du continu, $\beth_1$.

Étant donné l'\emph{indépendance} de la (G)CH
par rapport à $\ZFC$, que dire de sa
\emph{vérité} ? Il y aurait \emph{beaucoup}
à dire : de nombreux travaux récents et féconds
en théorie des ensembles portent précisément
sur ces questions. Deux remarques rapides
méritent néanmoins d'être soulignées.

Premièrement, ces résultats formels
d'indépendance n'impliquent \emph{pas
immédiatement} que la vérité de la GCH
ou de la CH soit \emph{indéterminée}.
Peut-être faut-il seulement ajouter des
axiomes qui nous paraissent naturels et
qui trancheront la question dans un sens
ou dans l'autre. Gödel lui-même pensait
que c'était la bonne réponse.

Deuxièmement, l'indépendance de la CH
par rapport à $\ZFC$ est assurément
\emph{frappante}, mais elle n'a rien
d'\emph{incroyable}, au sens littéral.
Simplement, d'après tout ce que nous dit
$\ZFC$, passer d'un cardinal à son
successeur peut demander un outil plus
fin que la seule opération qui forme
l'ensemble des parties.
% L'opération de formation de l'ensemble des parties
% fait passer d'une étape de la hiérarchie à la suivante
% (de $V_{\alpha}$ à $V_{\alpha+1}$).

Ces deux remarques faites, pour en savoir
plus, il vous faudra vous tourner vers les
philosophes et les mathématiciens engagés
dans ce débat.\footnote{Vous pourriez
commencer par lire \citet[\S15.6]{Potter2004}.}

\end{document}
